\section{The beta lattice and its truncations}\label{tr:sec:lattice}

Write $\xi_4(s)=\xi(s,\chi_4)=(4/\pi)^{(s+1)/2}\Gamma(\tfrac{s+1}2)L(s,\chi_4)$ for the completed Dirichlet beta function of Section~\ref{sec:dbeta:beta}, entire and symmetric about $s=\tfrac12$, and let $\omega$ be the exit time variable of Remark~\ref{rem:dbeta:exit} divided by $2\pi$. Then
\[
\E e^{-\lambda\omega}=\sech\sqrt{\lambda/2}=\prod_{k\ge1}\Big(1+\frac\lambda{\lambda_k}\Big)^{-1},\qquad\lambda_k=\lambda_*\big(k-\tfrac12\big)^2,\quad\lambda_*=2\pi^2,
\]
so that $\omega=\sum_ke_k/\lambda_k$ with independent standard exponential variables $e_k$, and the Thorin measure of $\omega$ is $\sum_k\delta_{\lambda_k}$. By Remark~\ref{rem:dbeta:exit} and the functional equation of $\xi_4$ the Mellin transform is
\begin{equation}\label{tr:eq:M}
M(z)=\E\,\omega^{-z}=(2\pi)^z\xi_4(2z),
\end{equation}
The nontrivial zeros of $L(s,\chi_4)$ are the points $z=\rho/2$ of the strip $0<\Re z<\frac12$ at which $M$ vanishes, and $g(z)=(2\pi)^{-z}M(z)=\xi_4(2z)$ is symmetric about $\Re z=\frac14$. GRH for $\chi_4$ says that they all lie on $\Re z=\frac14$. This is the lattice of Section~\ref{ch:rz4} with the atoms $\lambda_k$ in place of the squares $k^2$.

Split $\omega=\omega_n+T_n$ into the first $n$ atoms and the tail $T_n=\sum_{k>n}e_k/\lambda_k$. With $q_j(n)=\sum_{k>n}\lambda_k^{-j}$ the cumulants of the tail are $\kappa_j(T_n)=(j-1)!\,q_j(n)$, and by the midpoint rule (Lemma~\ref{tr:lem:tail})
\begin{equation}\label{tr:eq:cum}
\kappa_j(T_n)=\frac{(j-1)!}{\lambda_*^{\,j}(2j-1)\,n^{2j-1}}\Big(1-\frac{j(2j-1)}{12\,n^2}+O(n^{-4})\Big).
\end{equation}
A truncation replaces $T_n$ by a variable $R_n\ge0$ independent of $\omega_n$ and puts $M_n(z)=\E(\omega_n+R_n)^{-z}$. Three choices are used, the third a family.

\paragraph{Hard cutoff} $R_n=0$. With $A_k^{(n)}=\prod_{j\le n,\,j\ne k}\lambda_j/(\lambda_j-\lambda_k)$ the density of $\omega_n$ is $\sum_{k\le n}A_k^{(n)}\lambda_ke^{-\lambda_kw}$, and
\[
M_n^{(0)}(z)=\E\omega_n^{-z}=\Gamma(1-z)\sum_{k\le n}A_k^{(n)}\lambda_k^{\,z},
\]
an exponential polynomial times $\Gamma(1-z)$. It is holomorphic for $\Re z<n$ (the density behaves like $w^{n-1}$ at $0$); the poles of $\Gamma(1-z)$ at $1,\dots,n-1$ cancel and the pole at $n$ does not, so $M_n^{(0)}$ is not entire.

\paragraph{Mean correction} $R_n=c_n=\E T_n=q_1(n)=\psi'(n+\frac12)/(2\pi^2)$, a constant. Since $\int_0^\infty(w+c)^{-z}\lambda e^{-\lambda w}dw=\lambda^ze^{\lambda c}\Gamma(1-z,\lambda c)$,
\[
M_n^{(1)}(z)=\E(\omega_n+c_n)^{-z}=\sum_{k\le n}A_k^{(n)}\lambda_k^{\,z}e^{\lambda_kc_n}\,\Gamma(1-z,\lambda_kc_n),
\]
a finite sum of upper incomplete gamma functions, entire in $z$, of the same type as Haglund's approximants of $\Xi$ \cite{Haglund09} (Chapters~\ref{chap:miyatake} and \ref{chap:evt}) but with the weights of a Thorin lattice of $L(s,\chi_4)$.

\paragraph{Matching} $R_n\ge0$ independent of $\omega_n$ whose first $m$ cumulants equal those of $T_n$. For $m=0$ the hard cutoff is admissible, and $c_n$ is the deterministic choice for $m=1$; a mean match alone does not force $R_n$ to be constant. The distinction matters for $m\ge2$ (Section~\ref{tr:sec:drift}).

At $n=40$ and $z=0.3$ the mean corrected value is $M_{40}^{(1)}(0.3)=1.70347125$ against $M(0.3)=1.70347184$, and the hard cutoff at $n=200$ gives $M_{200}^{(0)}(0.3)=1.70461$ (computed, by the Laplace integral of the density with no cancellation of the $A_k^{(n)}$; two runs at different precision and step agree to $10^{-34}$; more digits: $1.703471255$, $1.70347184$, $1.704611458$).

\begin{theorem}[Local criterion]\label{tr:thm:local}
Let $R_n\ge0$ be independent of $\omega_n$ with $R_n\to0$ in probability, as for the hard cutoff and for every nonnegative remainder with $\E R_n=c_n$ (Markov's inequality gives $R_n\to0$ in probability, since $c_n=O(n^{-1})$). Then $M_n\to M$ locally uniformly on $0<\Re z<\frac12$, and GRH for $\chi_4$ holds if and only if every compact subset of $0<\Re z<\frac14$ is free of zeros of $M_n$ for all sufficiently large $n$.
\end{theorem}

\begin{proof}
Let $X_n=\omega_n+R_n$. Coupling the exponential summands gives $\omega_n\uparrow\omega$ almost surely, so $X_n\to\omega$ in probability, and the laws of $X_n$ converge weakly to that of $\omega$. Let $K$ be a compact subset of the strip and choose $b\in(\max_K\Re z,2)$. For $n\ge2$, $X_n\ge\omega_2$, so $|X_n^{-z}|\le1+\omega_2^{-b}$, which is integrable because the density of $\omega_2$ is $O(w)$ at $0$. Weak convergence and this uniform integrability give $\E X_n^{-z}\to\E\omega^{-z}$ for every $z\in K$. The same bound on a slightly larger compact set gives holomorphy and local boundedness, and Vitali's theorem gives local uniform convergence.

If GRH holds, $M$ has no zero on a compact $K\subset\{0<\Re z<\frac14\}$, so $\min_K|M|>0$ and $M_n\ne0$ on $K$ for large $n$. Conversely, if $M$ has a zero in $0<\Re z<\frac14$, a small closed disc about it with zero free boundary lies in the half strip, and Rouch\'e's theorem gives zeros of $M_n$ in the disc for every large $n$, contradicting the property ($M$ is not identically zero, being positive on the real segment). A zero in $\frac14<\Re z<\frac12$ is carried to the left half strip by the functional equation $(2\pi)^{-z}M(z)=(2\pi)^{z-\frac12}M(\frac12-z)$.
\end{proof}

The theorem, the analogue of Theorem~\ref{rz:thm:hurwitzrh} for $\chi_4$, turns GRH into a statement about explicit finite objects, but as an equivalence: the difficulty is to control the zeros of $M_n$ uniformly.

\section{The drift of the truncated zeros}\label{tr:sec:drift}

For $m=0$ and zeta the drift is Proposition~\ref{rz:prop:firstorder}. For $\chi_4$ and every order $m$ one point matters: a remainder matching the first $m$ cumulants of the tail has a first unmatched cumulant of its own, and the drift depends on the difference of the two.

\begin{lemma}\label{tr:lem:tail}
For fixed $j\ge1$, $q_j(n)=\lambda_*^{-j}\big\{n^{1-2j}/(2j-1)-\frac j{12}n^{-2j-1}+O(n^{-2j-3})\big\}$. In particular
\begin{align*}
c_n&=\frac1{\lambda_*n}-\frac1{12\lambda_*n^3}+O(n^{-5}),\\
v_n=\operatorname{Var}T_n=q_2(n)&=\frac1{3\lambda_*^2n^3}-\frac1{6\lambda_*^2n^5}+O(n^{-7}),\\
\kappa_3(T_n)&=\frac2{5\lambda_*^3n^5}+O(n^{-7}).
\end{align*}
\end{lemma}

\begin{proof}
For $p>1$, expand $x^{-p}$ about the midpoint of $[k-1,k]$ to fourth order: the quadratic term contributes $p(p+1)(k-\frac12)^{-p-2}/24$ and the remainder sums to $O(n^{-p-3})$. Summing over $k>n$ gives $\sum_{k>n}(k-\tfrac12)^{-p}=n^{1-p}/(p-1)-(p/24)n^{-p-1}+O(n^{-p-3})$, using the same estimate without the correction for $\sum_{k>n}(k-\frac12)^{-p-2}$. Take $p=2j$.
\end{proof}

Fix $m\ge0$ and write $r=m+1$, $p=2r-1$. Let $R_n\ge0$ be independent of $\omega_n$ with $\kappa_j(R_n)=\kappa_j(T_n)$ for $1\le j<r$, and assume, for a constant $C$ independent of $n$ and $j$,
\begin{equation}\label{tr:eq:A}
\log\E e^{-tR_n}=\sum_{j\ge1}\frac{(-1)^j\kappa_j(R_n)}{j!}t^j\quad(0\le t<n^2/C),\qquad\frac{|\kappa_j(R_n)|}{j!}\le C^jn^{1-2j}.
\end{equation}
This sufficient condition does not follow from matching finitely many cumulants. It holds for the hard cutoff, the constant remainder and the gamma remainders below, and for the tail $T_n$ itself after increasing $C$, since $\log\E e^{-tT_n}=-\sum_{k>n}\log(1+t/\lambda_k)$ and $q_j(n)\le(2j-1)^{-1}\lambda_*^{-j}(n-\frac12)^{1-2j}$. Put $\Delta_{r,n}=\kappa_r(T_n)-\kappa_r(R_n)$.

\begin{theorem}[Drift law]\label{tr:thm:drift}
Under these hypotheses, uniformly for $z$ in a compact subset $K$ of $\Re z>0$,
\begin{equation}\label{tr:eq:Mn}
M_n(z)-M(z)=-\frac{(-1)^r\Delta_{r,n}}{r!}\,(z)_r\,M(z+r)+O_K\big(n^{-p-2}+n^{-2p}\big),
\end{equation}
with $(z)_r=z(z+1)\cdots(z+r-1)$.
Let $\rho=\frac12+i\gamma$ be a simple zero of $\xi_4$ and $z_0=\rho/2$. For all large $n$ a fixed small disc about $z_0$ contains exactly one zero $z_n$ of $M_n$, it is simple, and
\begin{equation}\label{tr:eq:zn}
\begin{gathered}z_n-z_0=\frac{(-1)^r\Delta_{r,n}}{r!}(z_0)_r\frac{M(z_0+r)}{M'(z_0)}+O\big(n^{-p-2}+n^{-2p}\big),\\
\frac{M(z_0+r)}{M'(z_0)}=\frac{(2\pi)^r\xi_4(\rho+2r)}{2\,\xi_4'(\rho)}.
\end{gathered}\end{equation}
If $\Delta_{r,n}=d\,n^{-p}+o(n^{-p})$ with $d\ne0$, then $z_n-z_0=D_rn^{-p}+o(n^{-p})$ with $D_r=\dfrac{(-1)^rd\,(2\pi)^r(z_0)_r\,\xi_4(\rho+2r)}{2\,r!\,\xi_4'(\rho)}\ne0$.
\end{theorem}

\begin{proof}
The Laplace transforms satisfy $\E e^{-t(\omega_n+R_n)}=L(t)e^{D_n(t)}$ with $L(t)=\sech\sqrt{t/2}$ and $D_n=\log\E e^{-tR_n}-\log\E e^{-tT_n}$. The first $r-1$ cumulants cancel, so \eqref{tr:eq:A} gives, for $0\le t\le n^2/(2C)$, $D_n(t)=-\frac{(-1)^r\Delta_{r,n}}{r!}t^r+O(n^{-p-2}t^{r+1})$ with $\Delta_{r,n}=O(n^{-p})$. On $0\le t\le s_n=A(\log n)^2$ this is $o(1)$, and the exponential gives $e^{D_n}-1=-\frac{(-1)^r\Delta_{r,n}}{r!}t^r+O(n^{-p-2}t^{r+1}+n^{-2p}t^{2r})$. For $t\ge s_n$: with $B=\max_K\Re z$, an integer $Q>B$ and $J_n=\lfloor\sqrt{s_n/\lambda_*}\rfloor$, the first $Q$ factors of $\E e^{-t\omega_n}$ are at most $C_Qt^{-Q}$ and each of the factors with $Q<k\le J_n$ is at most $\frac12$, so $\E e^{-t\omega_n}\le C_Qt^{-Q}2^{-(J_n-Q)}$. Since $R_n\ge0$, both $L(t)$ and $\E e^{-t(\omega_n+R_n)}$ are bounded by it, and the Mellin integrals over $[s_n,\infty)$ are $O(n^{-N})$ for every $N$ once $A$ is large. The identity $\E X^{-z}=\Gamma(z)^{-1}\int_0^\infty t^{z-1}\E e^{-tX}dt$ converges absolutely for $\Re z>0$ and large $n$, and $\Gamma(z+r)/\Gamma(z)=(z)_r$, so integrating the expansion gives \eqref{tr:eq:Mn}.

By Rouch\'e's theorem a small disc about the simple zero $z_0$ contains exactly one zero $z_n$ of $M_n$, and it is simple. Put $E_n=M_n-M$; then $E_n=O(n^{-p})$ and, by Cauchy's estimate on a slightly larger disc, $E_n'=O(n^{-p})$. Thus $z_n-z_0=O(n^{-p})$, and Taylor's formula $0=M'(z_0)(z_n-z_0)+E_n(z_0)+O(n^{-2p})$ gives \eqref{tr:eq:zn}. The ratio follows from \eqref{tr:eq:M}: $M'(z_0)=2(2\pi)^{z_0}\xi_4'(\rho)$ and $M(z_0+r)=(2\pi)^{z_0+r}\xi_4(\rho+2r)$. Finally $(z_0)_r\ne0$, $\xi_4'(\rho)\ne0$, and $\xi_4(\rho+2r)\ne0$ because $\Re(\rho+2r)>1$, where the Euler product of $L(s,\chi_4)$ has no zero, so $D_r\ne0$ when $d\ne0$.
\end{proof}

\begin{corollary}[The two named schemes]\label{tr:cor:named}
Let $\rho=\frac12+i\gamma$ be a simple zero of $\xi_4$. Then
\begin{equation}\label{tr:eq:delta}
\delta_m=\frac{(-1)^{m+1}(z_0)_{m+1}\,\xi_4(\rho+2m+2)}{2(m+1)(2m+1)\,\pi^{m+1}\,\xi_4'(\rho)},
\end{equation}
and for the hard cutoff $z_n=z_0+\delta_0/n+O(n^{-2})$, for the mean correction $z_n=z_0+\delta_1/n^3+O(n^{-5})$. For a remainder that satisfies the hypotheses of Theorem~\ref{tr:thm:drift} at order $m$ and has $\kappa_{m+1}(R_n)=o(n^{-(2m+1)})$ the displacement is $\delta_m n^{-(2m+1)}(1+o(1))$.
\end{corollary}

\begin{proof}
If $\kappa_r(R_n)=o(n^{-p})$, then $d=(r-1)!/[\lambda_*^r(2r-1)]$ by Lemma~\ref{tr:lem:tail}, and $D_r$ of Theorem~\ref{tr:thm:drift} is $\delta_m$ of \eqref{tr:eq:delta}. For the hard cutoff $r=1$ and $\Delta_{1,n}=c_n$; for the constant $c_n$, $r=2$, $\Delta_{2,n}=v_n$. The error is $O(n^{-3}+n^{-2})$ for $r=1$ and $O(n^{-5}+n^{-6})$ for $r=2$. For the hard cutoff the relative error is $O(n^{-1})$, not $O(n^{-2})$: the error term $n^{-2p}=n^{-2}$ of \eqref{tr:eq:Mn} is the term $\frac12c_n^2(z)_2M(z+2)$ of the expansion of $e^{D_n}$, which is not cancelled by the variance term $v_n=O(n^{-3})$ and has a nonzero coefficient at $z_0$.
\end{proof}

The first zero of $L(s,\chi_4)$ is at $\gamma=6.0209489$. Table~\ref{tr:tab:drift} compares $\delta_m$ with the measured $n^{2m+1}(z_n-z_0)$ (computed, precision 40 digits with step $h=1/4$ against 50 digits with step $h=1/8$, differences below $10^{-26}$). The hard cutoff approaches $\delta_0$ with relative error $(0.27+0.51i)/n$, of modulus $0.58/n$, and the mean correction with an error that falls by a factor of $4$ at each doubling of $n$, as the orders $O(n^{-1})$ and $O(n^{-2})$ say; the finite $n$ values agree with the predicted constants to two or three digits, not four. For the hard cutoff the real part of $z_n-\frac14$ is $9.44\times10^{-3}$ at $n=128$ and $4.71\times10^{-3}$ at $n=256$. For the mean correction it is $5.87\times10^{-5}$, $3.10\times10^{-4}$, $4.72\times10^{-5}$, $6.15\times10^{-6}$, $7.76\times10^{-7}$, $9.73\times10^{-8}$ at $n=4,8,16,32,64,128$. The value at $n=4$ is small by accident (the law $\Re\delta_1/n^3$ gives $3.2\times10^{-3}$), and the sequence is not monotone. The hard cutoff would need $n\approx\Re\delta_0/(8\times10^{-7})\approx1.5\times10^6$ to reach $8\times10^{-7}$ (from the asymptotic law, not a run).

\begin{table}[h]\centering\small
\fitwidth{\begin{tabular}{rccc}
\toprule
$n$ & hard cutoff, $n(z_n-z_0)$ & mean correction, $n^3(z_n-z_0)$ & gamma remainder, $n^5(z_n-z_0)$\\
\midrule
$16$ & $1.2424-0.5466i$ & $0.1934-1.0874i$ & $-1.1423-0.4084i$\\
$32$ & $1.2234-0.5618i$ & $0.2015-1.0853i$ & $-1.1429-0.4370i$\\
$64$ & $1.2137-0.5692i$ & $0.2035-1.0849i$ & $-1.1428-0.4443i$\\
$128$ & $1.2089-0.5728i$ & $0.2040-1.0848i$ & $-1.1428-0.4461i$\\
\midrule
limit & $\delta_0=1.2041-0.5763i$ & $\delta_1=0.2041-1.0848i$ & $\tfrac49\delta_2=-1.1428-0.4467i$\\
\bottomrule
\end{tabular}}
\caption{The drift of the first zero of $M_n$ at $\gamma=6.0209489$ (computed). The last column is the gamma remainder of Proposition~\ref{tr:prop:gamma}, for which $\delta_2=-2.5713-1.0051i$ is not the limit.}\label{tr:tab:drift}
\end{table}

\subsection{Why the matching hypothesis cannot be dropped}

The note \cite{P26trunc} states \eqref{tr:eq:delta} for every $m$ and every remainder matching $m$ cumulants, for instance a generalised gamma convolution with a few atoms. This is false for $m\ge2$, because the first unmatched cumulant of $R_n$ enters.

\begin{proposition}\label{tr:prop:gamma}
Let $R_n$ have the gamma law with shape $c_n^2/v_n$ and scale $v_n/c_n$, so that $\E R_n=c_n$ and $\operatorname{Var}R_n=v_n$ exactly. It satisfies \eqref{tr:eq:A}, and with $r=3$ the drift law gives $z_n=z_0+\frac49\delta_2\,n^{-5}+O(n^{-7})$. A second remainder $R_n=G_{1,n}+G_{2,n}$, a sum of two independent gamma variables with matched mean and variance and the same natural scale, gives $z_n=z_0-\frac19\delta_2\,n^{-5}+O(n^{-7})$. A bounded remainder $R_n=b_n+hB_n$ with $B_n$ a Bernoulli variable, $h>0$ small and fixed, matched to the mean and variance, has $z_n-z_0=cn^{-3}+O(n^{-5})$ with $c\ne0$.
\end{proposition}

\begin{proof}
For the gamma law $\kappa_3(R_n)=2v_n^2/c_n$, so $\Delta_{3,n}=2\big(q_3-q_2^2/q_1\big)=\frac8{45\lambda_*^3n^5}+O(n^{-7})=(\frac49+O(n^{-2}))\kappa_3(T_n)$ by Lemma~\ref{tr:lem:tail}, and the difference in brackets is positive for every $n$ by the Cauchy Schwarz inequality applied to $\lambda_k^{-1/2}$ and $\lambda_k^{-3/2}$. Theorem~\ref{tr:thm:drift} gives the first claim. For two gamma variables with shapes and scales $\big(\frac{8c_n}{5s_n},\frac{s_n}2\big)$ and $\big(\frac{c_n}{15s_n},3s_n\big)$, $s_n=v_n/c_n$, the mean is $c_n$, the variance is $v_n$ and $\kappa_3(R_n)=4v_n^2/c_n$, so $\Delta_{3,n}=-\frac2{45\lambda_*^3n^5}+O(n^{-7})=(-\frac19+O(n^{-2}))\kappa_3(T_n)$. For the Bernoulli remainder put $p_n=\frac12(1-\sqrt{1-4v_n/h^2})$ and $b_n=c_n-hp_n$, so that $R_n\ge0$ for large $n$ with mean $c_n$ and variance $v_n$ exactly; then $\kappa_3(R_n)=hv_n(1-2p_n)$ is of order $n^{-3}$. Write $M_h(z)=\E(\omega+h)^{-z}$ and $H_h(z)=h^{-2}\big[M_h(z)-M(z)+hzM(z+1)\big]-\frac12(z)_2M(z+2)$. The method of the proof of Theorem~\ref{tr:thm:drift} gives $M_n(z)-M(z)=v_nH_h(z)+O_K(n^{-5})$, and Taylor's theorem gives $H_h(z)=-\frac h6(z)_3M(z+3)+O(h^2)$, which does not vanish at $z_0$ for small $h$; hence $z_n-z_0=-H_h(z_0)/(3\lambda_*^2M'(z_0))\,n^{-3}+O(n^{-5})$ with a nonzero coefficient.
\end{proof}

The gamma column of Table~\ref{tr:tab:drift} agrees: $n^5(z_n-z_0)/\delta_2=0.4444$ at $n=128$, against $\frac49$ (computed). No generalised gamma convolution repairs the formula for $m=2$: for every positive measure $U$ the moments $m_j=\int u^{-j}U(du)$ satisfy $m_1m_3\ge m_2^2$, so a generalised gamma convolution remainder with the mean and variance of $T_n$ has $\kappa_3(R_n)\ge2v_n^2/c_n$, of exact order $n^{-5}$ (proved, by the Cauchy Schwarz inequality), with the one gamma remainder as the extreme case. So $\delta_m$ for $m\ge2$ belongs to a hypothetical remainder with $\kappa_{m+1}(R_n)=o(n^{-(2m+1)})$, not to the generalised gamma convolutions the note proposes.

\section{The side of approach is a Gram parity}\label{tr:sec:phase}

For zeta, Theorem~\ref{rz:thm:kappaexact} makes the sign of $\Re\kappa(\rho)$, up to the phase of $\zeta(\frac52+i\gamma)$, that of $\sin\theta(\gamma)/Z'(\gamma)$, a Gram parity. The same holds for $\chi_4$ and every $m$, with one more correction. Put $G(s)=(4/\pi)^{s/2}\Gamma(\frac{s+1}2)$, so that $\xi_4=\frac2{\sqrt\pi}GL$, and let $\vartheta_\beta(x)=\arg G(\tfrac12+ix)$ and $Z_\beta(x)=e^{i\vartheta_\beta(x)}L(\tfrac12+ix,\chi_4)$ be the Riemann Siegel and Hardy functions of $\chi_4$ ($\vartheta$ and $Z$ are those of zeta in Section~\ref{sec:ch5:turing}). Then $G(\frac12+ix)L(\frac12+ix,\chi_4)=|G|\,Z_\beta(x)$ is real, and the Gram points of $\chi_4$ are the solutions of $\vartheta_\beta(x)=j\pi$.

\begin{theorem}[Phase law]\label{tr:thm:phase}
Let $\rho=\frac12+i\gamma$, $\gamma>0$, be a simple zero of $\xi_4$ and let $\delta_m$ be the number \eqref{tr:eq:delta}. Then
\[
\delta_m=\frac{i\,C_m(\gamma)\,e^{i\Theta_m(\gamma)}}{Z_\beta'(\gamma)},\qquad C_m(\gamma)>0,
\]
\[
\Theta_m(\gamma)=\vartheta_\beta(\gamma)+\arg L(\rho+2m+2,\chi_4)-\sum_{j=0}^m\Big[\arctan\frac{\frac12+2j}\gamma+\arctan\frac{\frac32+2j}\gamma\Big],
\]
so that $\operatorname{sgn}\Re\delta_m=\operatorname{sgn}\big[-\sin\Theta_m(\gamma)\,Z_\beta'(\gamma)\big]$, zero included. For an actual remainder satisfying the hypotheses of Theorem~\ref{tr:thm:drift} and with $\eta_m=\lim_n\big[\kappa_{m+1}(T_n)-\kappa_{m+1}(R_n)\big]/\kappa_{m+1}(T_n)$ the displacement is $\eta_m\delta_mn^{-(2m+1)}(1+o(1))$ when the limit exists and $\eta_m\ne0$, and the sign of the real part is $\operatorname{sgn}\big[-\eta_m\sin\Theta_m(\gamma)Z_\beta'(\gamma)\big]$ when $\eta_m\sin\Theta_m\ne0$.
\end{theorem}

\begin{proof}
The gamma factor satisfies $G(s+2)=\frac4\pi\cdot\frac{s+1}2G(s)$, so $\xi_4(\rho+2m+2)=\frac2{\sqrt\pi}(\frac4\pi)^{m+1}\prod_{j=0}^m\big(\frac{\rho+1}2+j\big)G(\rho)L(\rho+2m+2,\chi_4)$. On the critical line $G(\rho)=|G(\rho)|e^{i\vartheta_\beta(\gamma)}$, and differentiating $G(\frac12+ix)L(\frac12+ix,\chi_4)=|G|Z_\beta(x)$ at a zero gives $i\,GL'(\rho)=|G|Z_\beta'(\gamma)$, that is $\frac{\sqrt\pi}2\xi_4'(\rho)=-i|G(\rho)|Z_\beta'(\gamma)$. In \eqref{tr:eq:delta} the product $(z_0)_{m+1}\prod_j(\frac{\rho+1}2+j)=\prod_{j=0}^m(\frac\rho2+j)(\frac{\rho+1}2+j)$ has factors $\frac14+j+\frac{i\gamma}2$ and $\frac34+j+\frac{i\gamma}2$ with arguments $\frac\pi2-\arctan\frac{1/2+2j}\gamma$ and $\frac\pi2-\arctan\frac{3/2+2j}\gamma$ for $\gamma>0$. The terms $\frac\pi2$ total $(m+1)\pi$ and cancel $(-1)^{m+1}$, which leaves $\delta_m=iC_m e^{i\Theta_m}/Z_\beta'(\gamma)$ with $C_m=(4/\pi)^{m+1}|\prod|\,|L(\rho+2m+2,\chi_4)|/[2(m+1)(2m+1)\pi^{m+1}]$. Taking real parts, $\Re(ie^{i\Theta})=-\sin\Theta$. The last statement is Theorem~\ref{tr:thm:drift} with $d=\eta_m\lim_nn^{2m+1}\kappa_{m+1}(T_n)$.
\end{proof}

At two consecutive simple zeros $Z_\beta'$ has opposite signs, so the sign of the drift alternates with the parity of the zero index unless $\sin\Theta_m$ changes sign, which is the Gram law of $\chi_4$ read through $\Theta_m$. The two corrections to $\vartheta_\beta$ are the arithmetic phase $\arg L(\rho+2m+2,\chi_4)$ and the Pochhammer phase, which is of exact size $-2(m+1)^2/\gamma+O_m(\gamma^{-3})$ (proved, first order sums). The arithmetic phase is bounded, $|L(\sigma+i\gamma,\chi_4)-1|\le\epsilon_\sigma=\sum_{k\ge1}(2k+1)^{-\sigma}$, but it does not tend to zero at fixed $\sigma>1$: the mean square of $\Im L(\sigma+it,\chi_4)$ over $[0,T]$ tends to $\frac12\sum_{n\ge2}\chi_4(n)^2n^{-2\sigma}>0$ (proved, by orthogonality of distinct frequencies). For zeta the corresponding statement is that $\Re\kappa$ has the sign of $-\cos(\pi S_{\rm mid})$ up to a shift of the threshold $\frac12$ by about $0.11$ (Section~\ref{ch:rz4}).

\subsection{All zeros below height 500}

\begin{resultpar}\label{tr:res:gram}
The zeros of $L(s,\chi_4)$ below height $500$ were located by sign changes of $Z_\beta$ on two independent grids (steps $0.02$ and $0.0137$) and refined at two precisions (computed). There are $379$ zeros, all simple and on the line as far as the computation sees (the minimal gap is $0.1336$, against the step $0.02$), in agreement with $\vartheta_\beta(500)/\pi=379.15$ and with $N(500)=379.1$ from \eqref{eq:dbeta:count}. This is not a certified count: there is no Turing bound, and the certified count in this book stops at height $300$, where $N(300)=203$ (Section~\ref{gl:sec:numchi}). For each zero the numbers $\delta_m$, $0\le m\le5$, were computed from \eqref{tr:eq:delta}, and the sign law of Theorem~\ref{tr:thm:phase} holds at all $379$ zeros for each $m$ (an identity, so a check of the arithmetic).

Measuring the Gram parity by the index $\lfloor\vartheta_\beta(\gamma)/\pi\rfloor$ minus the zero number, $374$ zeros have the offset $-1$ and five do not (two have $-2$ and three have $0$). The five exceptions, at $\gamma=138.7502,290.2521,364.4197,400.8264,422.5104$, are listed in Table~\ref{tr:tab:gram}. For the mean correction ($m=1$), $\Re\delta_1<0$ exactly at these five zeros and at no other. For the hard cutoff ($m=0$), $\Re\delta_0<0$ at the last four, and at $138.750$ the arithmetic phase is large enough to turn the sign. For $m\ge2$ the Pochhammer phase produces further negative real parts at low height: $\Re\delta_m<0$ at $8$, $11$, $19$ and $29$ of the $379$ zeros for $m=2,3,4,5$ (for $m=2$ at $6.021$, $12.988$, $138.750$, $236.242$, $290.252$, $364.420$, $400.826$ and $422.510$, and for $m=5$ at $22$ zeros below height $105$). These counts concern the formal $\delta_m$; by Proposition~\ref{tr:prop:gamma} they give the drift of an actual remainder only after multiplication by $\eta_m$.\end{resultpar}

\begin{table}[h]\centering\small
\begin{tabular}{rrcccc}
\toprule
$\gamma$ & zero number & gap to the partner & $\Re\delta_0$ & $\Re\delta_1$\\
\midrule
$138.7502$ & $78$ & $0.2929$ & $>0$ & $<0$\\
$290.2521$ & $195$ & $0.4251$ & $<0$ & $<0$\\
$364.4197$ & $258$ & $0.3840$ & $<0$ & $<0$\\
$400.8264$ & $291$ & $0.1336$ & $<0$ & $<0$\\
$422.5104$ & $309$ & $0.1769$ & $<0$ & $<0$\\
\bottomrule
\end{tabular}
\caption{The five zeros of $L(s,\chi_4)$ below height $500$ with anomalous Gram parity (computed). The partner is the neighbouring zero at the stated gap (the left neighbour for the zero number $291$).}\label{tr:tab:gram}
\end{table}

Each of the five lies in a close pair of zeros (gaps recomputed). For zeta the analogue is the list of Section~\ref{ch:rz4}: three zeros among the first $460$ approached from the right, related to $|S_{\rm mid}|>0.56$.

\begin{corollary}[No one sided criterion]\label{tr:cor:nogo}
Fix $m\ge0$ and a remainder satisfying the hypotheses of Theorem~\ref{tr:thm:drift}, for which the limit $\eta_m$ exists. If some simple zero $\rho$ of $\xi_4$ on the critical line has $\Re(\eta_m\delta_m)<0$, then for all large $n$ the function $M_n$ has a zero in $0<\Re z<\frac14$, so no criterion of the form that $M_n$ has no zero in $0<\Re z<\frac14$ for infinitely many $n$ can hold at order $m$ for this remainder.
\end{corollary}

\begin{proof}
By Theorem~\ref{tr:thm:drift}, $\Re z_n=\frac14+n^{-(2m+1)}\Re(\eta_m\delta_m)+o(n^{-(2m+1)})$, which is below $\frac14$ for large $n$; and the zero lies in the strip $0<\Re z<\frac12$.
\end{proof}

This is the argument of Corollary~\ref{rz:cor:Tfalse} in the other direction. Its premise is computed, not proved: for the mean correction ($m=1$, $\eta_1=1$) the zero $\gamma=138.7502$ has $\Re\delta_1<0$ in floating point, and so do the four others. A zero with $\Re\delta_1<0$ proved in ball arithmetic would make the corollary unconditional for $m=1$, exactly as in Corollary~\ref{rz:cor:Tfalse}. The note states more, that infinitely many Gram anomalies of $\chi_4$ would force infinitely many left approaches at every fixed order $m$. This does not follow from the Gram parity alone, because the side is set by $\sin\Theta_m$ and not by $\sin\vartheta_\beta$: at $138.750$ the hard cutoff approaches from the right although the zero is anomalous. A precise implication is the following.

\begin{proposition}\label{tr:prop:gram}
Let $\sigma=2m+\frac52$, $\epsilon_\sigma=\sum_{k\ge1}(2k+1)^{-\sigma}<1$ and $B_m(\gamma)=\arcsin\epsilon_\sigma+\sum_{j=0}^m\big[\arctan\frac{1/2+2j}\gamma+\arctan\frac{3/2+2j}\gamma\big]$. Then $|\Theta_m(\gamma)-\vartheta_\beta(\gamma)|\le B_m(\gamma)$. If infinitely many simple critical zeros satisfy $-\sin\vartheta_\beta(\gamma)Z_\beta'(\gamma)<0$ and $\operatorname{dist}(\vartheta_\beta(\gamma),\pi\mathbb Z)>B_m(\gamma)$, then infinitely many zeros have $\Re\delta_m<0$.
\end{proposition}

\begin{proof}
The disc of centre $1$ and radius $\epsilon_\sigma$ contains $L(\sigma+i\gamma,\chi_4)$, and it is seen from the origin under the angle $\arcsin\epsilon_\sigma$, which bounds $|\arg L|$. The interval between $\vartheta_\beta(\gamma)$ and $\Theta_m(\gamma)$ then contains no multiple of $\pi$, the two sines have the same nonzero sign, and Theorem~\ref{tr:thm:phase} applies.
\end{proof}

The hypotheses of Proposition~\ref{tr:prop:gram} on the five ordinates, with the margin $B_m(\gamma)$, are not checked here. Whether Gram anomalies with the margin occur infinitely often for $\chi_4$ is open. For zeta Gram's law fails with positive frequency \cite{Trudgian11b}, and the analogy with the five failures below $500$ is a reason to expect the same for $\chi_4$ (heuristic). The comparison with the de Bruijn Newman constant, that these are the configurations that make $\Lambda\ge0$, is an analogy, and the theorem of Rodgers and Tao \cite{RodgersTao} does not prove a truncation obstruction.

\section{Transient zeros}\label{tr:sec:transient}

Besides the zeros attached to those of $L$, the truncations have zeros that wander into the left half strip. The counts of Table~\ref{tr:tab:counts} are zeros of $M_n$ in the rectangles $0.02<\Re z<0.2499$ (left) and $0.2501<\Re z<0.48$ (right) by the argument principle (computed, enclosures of the function values in ball arithmetic but a floating point winding sum, a second run at $1.5$ times the precision agrees; not certified).

\begin{table}[h]\centering\small
\begin{tabular}{lcccc}
\toprule
truncation & $n$ & $\Im z$ & left & right\\
\midrule
hard cutoff & 32 & $15$ to $60$ & 2 & 14\\
hard cutoff & 64 & $15$ to $60$ & 2 & 15\\
mean correction & 32 & $0.5$ to $20$ & 0 & 13\\
mean correction & 16 & $20$ to $40$ & 3 & 3\\
mean correction & 32 & $20$ to $40$ & 1 & 20\\
mean correction & 64 & $20$ to $40$ & 0 & 23\\
mean correction & 64 & $40$ to $75$ & 4 & 37\\
mean correction & 128 & $20$ to $40$ & 0 & 23\\
\bottomrule
\end{tabular}
\caption{Zeros of $M_n$ in the left and right half strips. The counts are computed at two precisions with identical counts. The source note printed $22$ for the hard cutoff at $n=32$, a blank at $n=64$ and $14$ for the mean correction at $n=32$, $\Im z$ from $0.5$ to $20$.}\label{tr:tab:counts}
\end{table}

The zeros of $L(s,\chi_4)$ with $0<\gamma<40$, $40<\gamma<80$ and $80<\gamma<150$ number $14$, $23$ and $48$ (computed, agreeing with \eqref{eq:dbeta:count}). The note's right counts for the mean correction reach $23$ in the band $20<\Im z<40$ at $n=64$ and $n=128$, and are $14$ in the first band at $n=32$, where the check of this chapter reported $13$ (the mean correction counts are computed, see the caption). A count of $13$ is not explained by a resolution frontier: the zero of $M_{32}^{(1)}$ near $\rho_1/2$ has $\Re z-\frac14=6.15\times10^{-6}$ (Table~\ref{tr:tab:drift}) and lies in the band $0.2499<\Re z<0.2501$ that the rectangles omit. At $n=64$ the band $40<\Im z<75$ holds $37$ captured zeros of the $48$ in $80<\gamma<150$ and four zeros in the left half strip (computed, Table~\ref{tr:tab:counts}). The note gives their ordinates (heights doubled) as $118.82$, $126.01$, $138.65$ and $142.92$, which were not recomputed (unchecked). Each is said to lie between the members of a close pair of zeros of $L$, namely $(118.538,119.453)$, $(125.769,126.299)$, $(138.457,138.750)$ and $(142.394,143.329)$; these eight ordinates were recomputed and are zeros of $L$. The third pair contains the first Gram anomaly. The pairs $(44.618,45.600)$ and $(67.637,68.366)$ are also pairs of zeros of $L$ (checked); the note reports that at $n=32$ the hard cutoff has transient zeros between their members, at $0.2356+22.5622i$ and $0.2411+34.0550i$, the first gone by $n=256$ and the second at $\Re z=0.237$ and moving right at $n=768$ (unchecked: this was not recomputed, since the hard cutoff at $n=768$ needs about $1500$ digits).

The note's picture is that mean correction reduces the transients of the hard cutoff and moves them to a resolution frontier that rises with $n$. In Table~\ref{tr:tab:counts} the left count in the band $20<\Im z<40$ is $3,1,0,0$ for $n=16,32,64,128$ (computed), consistent with a frontier that rises with $n$, while the hard cutoff keeps two left zeros at both $n=32$ and $n=64$ in $15<\Im z<60$. The rectangles of the count do not contain a neighbourhood of the critical line or of the left boundary. A zero $z_n=\frac14+\delta_1n^{-3}+o(n^{-3})$ with $\Re\delta_1<0$ eventually lies in $0.2499<\Re z<\frac14$ and is invisible to the count while remaining to the left of the line; so an empty left rectangle at a fixed height does not exclude such zeros, which are consistent with Theorem~\ref{tr:thm:local} because its compact sets stay away from the line. No uniform frontier is proved, and no statement that the transients occur only at close pairs.

\section{Symmetrised truncations}\label{tr:sec:sym}

Impose the functional equation on each truncation: put $g_n(z)=(2\pi)^{-z}M_n(z)=\E e^{-zV_n}$, with $V_n=\log(2\pi(\omega_n+R_n))$, and
\[
F_n(z)=g_n(z)+g_n(\tfrac12-z).
\]
Then $F_n(\frac12-z)=F_n(z)$ and $F_n(\bar z)=\overline{F_n(z)}$ on the common domain, so $F_n$ is real on $\Re z=\frac14$ and its zeros off the line come in quadruples; this is Proposition~\ref{rz:prop:symtrunc} for $\chi_4$. If $p_n$ is the density of $V_n$ and $q_n(v)=e^{-v/4}p_n(v)$, then $F_n(\frac14+iy)=\int_{\mathbb R}[q_n(v)+q_n(-v)]e^{-iyv}dv$, so $F_n(\frac14+iy)/F_n(\frac14)$ is the characteristic function of the positive even density $[q_n(v)+q_n(-v)]/F_n(\frac14)$ (proved, by the change of variables; for a remainder without a density use measures).

\begin{theorem}[Hurwitz criterion]\label{tr:thm:symhurwitz}
If for infinitely many $n$ all zeros of $F_n$ (for the remainders of Theorem~\ref{tr:thm:local}) in $0<\Re z<\frac12$ lie on $\Re z=\frac14$, then GRH holds for $\chi_4$.
\end{theorem}

\begin{proof}
By Theorem~\ref{tr:thm:local} and the functional equation $F_n\to2\xi_4(2z)$ locally uniformly in the strip. Along the subsequence apply Hurwitz's theorem separately to the two open half strips $0<\Re z<\frac14$ and $\frac14<\Re z<\frac12$: the nonzero limit $2\xi_4(2z)$ has no zero in either. Hurwitz's theorem is not a theorem on a closed line, so this uses the subsequence and does not apply to one fixed approximant.
\end{proof}

\begin{proposition}\label{tr:prop:symlocal}
Every simple critical zero of $\xi_4(2z)$ is, for all large $n$, the limit of exactly one zero of $F_n$ in a small disc about it, and that zero lies on $\Re z=\frac14$.
\end{proposition}

\begin{proof}
Rouch\'e's theorem in a disc centred on the line gives exactly one zero of $F_n$ in it, and the reflection $z\mapsto\frac12-\bar z$ preserves the disc and the zero set, so the zero lies on the line. This holds even when the corresponding zero of $g_n$ approaches from the left.
\end{proof}

The proposition is local: on a fixed compact region it makes every zero of $F_n$ critical for large $n$, but says nothing about heights growing with $n$.

For the hard cutoff, counting all zeros in $0.02<\Re z<0.48$, $0.5<\Im z<60$ against sign changes on the line gives Table~\ref{tr:tab:sym} (computed, the same at $1.5$ times the precision). The count $64$ equals the number of zeros of $L(s,\chi_4)$ below height $120$ found by the same procedure. For $n=8$ and $16$ every zero of $F_n$ in this rectangle is critical, and the zeros of $F_8$ already approximate those of $L$ (Table~\ref{tr:tab:F8}, computed, heights $\gamma=2\Im z$).

\begin{table}[h]\centering\small
\begin{tabular}{rccc}
\toprule
$n$ & zeros in the strip & on the line & left zeros of $g_n$\\
\midrule
8 & 64 & 64 & 0\\
16 & 64 & 64 & 0\\
32 & 64 & 60 & 2\\
64 & 64 & 60 & 2\\
\bottomrule
\end{tabular}
\caption{Symmetrised hard cutoff (computed).}\label{tr:tab:sym}
\end{table}

\begin{table}[h]\centering\small
\begin{tabular}{lcccccc}
\toprule
$F_8$ & 5.9438 & 10.3288 & 12.9300 & 16.5280 & 18.1556 & 21.5568\\
$L(s,\chi_4)$ & 6.0209 & 10.2443 & 12.9876 & 16.343 & 18.292 & 21.451\\
\bottomrule
\end{tabular}
\caption{The first zeros of $F_8$ and of $L(s,\chi_4)$ (computed).}\label{tr:tab:F8}
\end{table}

\begin{remark}[Zeros of $g_n$ to the left]\label{tr:rem:biehler}
For $n=32$ and $64$ the four off line zeros of $F_n$ and the two zeros of $g_n$ to the left of the line are computed counts. The note reads this as a mechanism, that a zero of $g_n$ to the left of the line produces an off line pair of $F_n$, in the manner of the Hermite Biehler theorem. That cannot be a general statement. At a zero $z_*$ of $g_n$ one has $F_n(z_*)=g_n(\frac12-z_*)$, which need not vanish; and the Hermite Biehler theorem gives real zeros of $E+E^*$ from the absence of zeros of $E$ in a half plane, not the converse. Here is a counterexample with a positive transform. Let $w=z-\frac14$, $c=e^{-1/8}$ and $g(z)=(e^w+c)/(e^{-1/4}+c)$, a positive two point Laplace transform with $g(0)=1$ whose zeros $z=\frac18+(2k+1)\pi i$ lie to the left of the line. Its symmetrisation is $F(z)=2(\cosh w+c)/(e^{-1/4}+c)$, and every solution of $\cosh w=-c$ with $0<c<1$ has $\Re w=0$ (if $w=x+iy$ with $x\ne0$ then $\sin y=0$ and $|\cosh w|\ge1$). So all zeros of $F$ are critical (proved). In the data the off line zeros of $F_{32}$ and $F_{64}$ are four, accompanying two left zeros of $g_n$; that they coincide is an observation of the note, not recomputed.
\end{remark}

The modulus inequality $|g_n(\frac12-z)|<|g_n(z)|$ on an open set off the line implies $F_n\ne0$ there, since $|F_n(z)|\ge|g_n(z)|-|g_n(\frac12-z)|>0$ (proved). It is sufficient and not necessary. The note reports that it fails near the real axis even for $n=8$ and gives $-0.118$ at $z=0.01+0.5i$ for the minimum of $\log|g_8(z)|-\log|g_8(\frac12-z)|$. The recomputed minimum on the grid $x\in[0.01,0.25]$, $y\in[0.5,60.5]$ (steps $0.01$ and $0.05$) is $-0.0282$ at the same point (computed): the inequality fails, but with the smaller value. Its failure does not imply an off line zero. The note adds that $F_8$ has no zeros there because $g_8$ is positive on the real axis; positivity on the real segment gives a zero free neighbourhood $|\Im z|<c/B$ of it only with a bound $|F_8'|\le B$ and $c=\min F_8$ on the segment, and neither is given.

With the mean correction, below height $40$ ($\gamma<80$) the function $F_{64}$ has all $37$ zeros on the critical line and $F_{32}$ has $35$ of $37$ (computed, at two precisions with identical counts, in agreement with the note). The note's reading, that the symmetrised mean corrected truncations satisfy the hypothesis below their resolution frontier, rests on two values of $n$ and is an observation.

\section{Where this leaves the route}\label{tr:sec:path}

Theorem~\ref{tr:thm:local} is the form of Proposition~\ref{ev:prop:reduction} and of Section~\ref{ev:sec:path} for $\chi_4$, with the Dirichlet beta lattice in place of the folded Gumbel kernels. Theorem~\ref{tr:thm:drift} gives the exact rate of convergence of the truncated zeros for the hard cutoff and the mean correction, and under the matching hypothesis for every order; Theorem~\ref{tr:thm:phase} identifies the side as a Gram parity, and to height $500$ the drift of the mean correction has negative real part exactly at the five anomalous zeros. What is not proved is that the premise of Corollary~\ref{tr:cor:nogo} holds (a zero with $\Re\delta_1<0$ in ball arithmetic) and that Gram anomalies with a margin recur. The symmetrised truncations avoid the one sided obstruction by Proposition~\ref{tr:prop:symlocal}, and Theorem~\ref{tr:thm:symhurwitz} reduces GRH for $\chi_4$ to a statement for infinitely many $n$ that is neither proved nor tested beyond the heights of Table~\ref{tr:tab:sym}. Nothing in the computation is certified.

\subsection{Open problems on the truncations}

\begin{problem}\label{tr:prob:gram}
Prove that Gram anomalies with the margin $B_m(\gamma)$ of Proposition~\ref{tr:prop:gram} occur infinitely often for $L(s,\chi_4)$, adapting the arguments known for $\zeta$ \cite{Trudgian11b}. Together with Proposition~\ref{tr:prop:gram} this would make Corollary~\ref{tr:cor:nogo} unconditional at each order $m$ with $\eta_m>0$.
\end{problem}

\begin{problem}\label{tr:prob:certify}
Prove in ball arithmetic that $\Re\delta_1<0$ at the zero $\gamma=138.7502$ of $L(s,\chi_4)$, as is done for zeta in Corollary~\ref{rz:cor:Tfalse}, and certify the count of $379$ zeros below height $500$ with a Turing bound.
\end{problem}

\begin{problem}\label{tr:prob:F8}
Prove that $F_8$, the symmetrised hard cutoff with eight atoms, has only critical zeros, by combining the phase of $g_8$ at low height with the modulus inequality and a Newton polygon bound at large height. The analogues for zeta are Theorem~\ref{rz:thm:Yn} and Corollary~\ref{rz:cor:Ssmall}.
\end{problem}

\begin{problem}\label{tr:prob:largest}
Determine the largest $n$ for which the symmetrised mean corrected truncation has only critical zeros in the strip, and decide whether the first failure is related to the first Gram anomaly at $138.75$.
\end{problem}

\begin{problem}\label{tr:prob:remainder}
Find a remainder $R_n$, not a cumulant match, whose drift $\eta\delta_m$ has a sign independent of the Gram parity. By Theorem~\ref{tr:thm:phase} this needs a correction whose phase cancels $\vartheta_\beta(\gamma)$, which suggests building $R_n$ from the functional equation and not from the tail.
\end{problem}

\begin{problem}\label{tr:prob:zeta}
Carry out the analysis of orders $m\ge1$ for zeta with the lattice of Section~\ref{ch:rz4}, where $m=0$ is Theorem~\ref{rz:thm:kappaexact}, and compare the Gram parity of the corrected phase with the classical failures of Gram's law for zeta (the first near the Gram point $g_{126}\approx282.455$, where $Z(g_{126})=-0.0276$).
\end{problem}
